%% =====================================================================
%% Tagged PDF (PDF/UA): what LaTeX's tagging code does not handle by itself (yet), and where its look
%% differs from the untagged editions. Read by template/packages.tex, right after listings is loaded; everything
%% here does something only when tagging is on (\DocumentMetadata{..., tagging=on}).
%%   1. Code listings
%%   2. Equation numbers in the MathML (with tagging.lua)
%%   3. Theorems and proofs straight after one another (LaTeX's code of June 2026)
%%   4. Figures side by side
%%   5. No "Rerun" warning for tagpdf's running totals
%%   6. Table cells: their scope and column span also on the cell itself (for PAC)
%% =====================================================================
\ExplSyntaxOn
\tag_if_active:T { \cs_gset:Npn \thesistagging {} }
\ExplSyntaxOff
\ifdefined\thesistagging
\makeatletter

%% ---------------------------------------------------------------------------------------------
%% 1. Code listings: a code block (a verbatim structure of code lines), the line numbers as labels and
%% the caption as the listing's caption. LaTeX brings this itself from November 2026 (latex-lab-listings,
%% loaded with listings); for the June 2026 release, the same code, from the LaTeX team's set-up for the
%% tagpdf documentation (tagpdfdocu-patches.sty, 2026-08-21, LaTeX Project Public License). Without it,
%% every listing closed the section around it in the tagged structure.
\IfPackageLoadedTF{latex-lab-listings}{\let\thesis@lablistings\relax}{}
\ifdefined\thesis@lablistings\else
\ExplSyntaxOn
\NewTaggingSocket{listings/number}{2}
\NewTaggingSocketPlug{listings/number}{default}
  {
    \tag_mc_end_push:
    \tagstructbegin{tag=Lbl}
    \tagmcbegin{}
    #2
    \tagmcend
    \tagstructend
    \tag_mc_begin_pop:n{}
  }
\AssignTaggingSocketPlug{listings/number}{default}
\ExplSyntaxOff
\DeclareInstance{blockenv}{lstlisting}{std}{%
    name                   = lstlisting,
    tag-name               = verbatim,
    tag-attr-class         = ,
    tagging-recipe         = standard,
    inner-level-counter    = ,
    transparent-level      = true,
    block-instance         = lstlisting,
    inner-instance         = ,
    final-code             = \AssignStructureRole{para/textblock}
                               {\UseStructureName{block/verbatim/codeline}}\ignorespaces,
    tagging-suppress-paras = true}
\tagpdfsetup{float/new=lstlisting}
\lst@Key{numbers}{none}{%
    \let\lst@PlaceNumber\@empty
    \lstKV@SwitchCases{#1}%
    {none:\\%
     left:\def\lst@PlaceNumber{\llap{\normalfont
                \UseTaggingSocket{listings/number}{}{\lst@numberstyle{\thelstnumber}}\kern\lst@numbersep}}\\%
     right:\def\lst@PlaceNumber{\rlap{\normalfont
                \kern\linewidth \kern\lst@numbersep
                \UseTaggingSocket{listings/number}{}{\lst@numberstyle{\thelstnumber}}}}%
    }{\PackageError{Listings}{value `#1' for keyword `numbers' not supported}\@ehc}}
\lst@AddToHook{PreInit}{\ifhmode \UseTaggingSocket{inline/begin}{tag=\UseStructureName{text/verb}}\fi}
\lst@AddToHook{DeInit}{\ifhmode \UseTaggingSocket{inline/end}%
    \else \edef\@current@float@struct{\UseName{tag_get:n}{struct_num}}\BlockEnvEnd \fi}
\def\lst@makecaption{\AssignStructureRole{para/textblock}{Span}\def\@captype{lstlisting}\@makecaption}
\def\lst@maketitle{\def\@captype{lstlisting}%
    \AssignStructureRole{para/textblock}{Span}%
    \AssignTaggingSocketPlug{caption/label/begin}{noop}%
    \AssignTaggingSocketPlug{caption/label/end}{noop}%
    \AssignSocketPlug{caption/label}{noop}%
    \@makecaption\lst@title@dropdelim}
\fi
%% The spacing, as without tagging (both releases): listings adds its own skips and break points
%% (\lstset aboveskip, belowskip), so the block around a listing adds none ...
\DeclareInstance{block}{lstlisting-1}{std}{%
    left-margin = 0pt, para-vspace = 0pt, begin-vspace = 0pt, end-vspace = 0pt,
    begin-extra-vspace = 0pt, end-extra-vspace = 0pt, begin-unchained-vspace = 0pt,
    begin-penalty = 10000, end-penalty = 10000}
\DeclareInstanceCopy{block}{lstlisting-2}{lstlisting-1}
\DeclareInstanceCopy{block}{lstlisting-3}{lstlisting-1}
\DeclareInstanceCopy{block}{lstlisting-4}{lstlisting-1}
\DeclareInstanceCopy{block}{lstlisting-5}{lstlisting-1}
\DeclareInstanceCopy{block}{lstlisting-6}{lstlisting-1}
%% ... no stretch above it (the block adds the paragraph skip), and no skips of nothing (they stayed at
%% the foot of a page that breaks at the listing)
\def\lst@EveryDisplay{%
    \MathCollectFalse
    \parskip\z@skip
    \ifx\lst@float\relax \SimpleBlockEnv{lstlisting}{}%
    \else \SimpleBlockEnv{lstlisting}{tagging-recipe = standalone}\fi
    \thesis@unskipzero}
\def\thesis@unskipzero{\let\thesis@next\relax
    \ifvmode\ifnum\lastnodetype=11 % (a skip)
        \thesis@ifzeroskip\lastskip{\unskip\let\thesis@next\thesis@unskipzero}{}\fi\fi
    \thesis@next}

%% Spaces in code as space characters (both releases): listings sets each word in its own box, so the
%% space between two words is not next to a letter and LaTeX's tagging adds no space character there.
%% The tagged text read "importnumpyasnp"; with \pdffakespace (a space character, no width) it reads
%% "import numpy as np". The pages do not change. (Indentation stays without characters.)
\def\lst@outputspace{\pdffakespace\ }

\ExplSyntaxOn
%% ---------------------------------------------------------------------------------------------
%% 2. Equation numbers in the MathML (only when the MathML is made: luamml, math/setup in
%% \DocumentMetadata, with the OpenType maths fonts).
%%   multline: luamml leaves multline out of the environments whose equation number gets its own
%%   structure, but still refers to that structure when it writes the MathML ("structure with label 0
%%   is unknown"): treat multline like equation and gather.
%%   equation with split: LaTeX keeps only the split's table in the MathML and drops the equation
%%   number; it is kept back here and added again around that table (tagging.lua).
%%   tagging.lua also makes a number one <mn> (siunitx's 6.022\,140\,76 was five pieces), a word in
%%   upright letters one <mi> (\mathrm{disc}), and leaves out boxes with nothing to read.
\AddToHook{begindocument}
  {
    \IfPackageLoadedT{luamml}
      {
        \tl_if_exist:cF { c__luamml_label_multline_tl }
          { \tl_const:cn { c__luamml_label_multline_tl } { } }
        \directlua{dofile(kpse.find_file('template/tagging.lua'))}
        \NewTaggingSocketPlug{math/luamml/mtable/tag/save}{thesis}
          {
            \bool_if:NTF \l__math_gathersplit_bool
              { \use:c{thesis@luamml@savetag} \l__luamml_amsmath_tag_struct_num_int }
              { \__luamml_amsmath_save_tag_with_struct_elem:N \l__luamml_amsmath_tag_struct_num_int }
          }
        \AssignTaggingSocketPlug{math/luamml/mtable/tag/save}{thesis}
        \NewTaggingSocketPlug{math/luamml/mtable/tag/set}{thesis}
          { \bool_if:NF \l__math_gathersplit_bool { \__luamml_amsmath_set_tag: } }
        \AssignTaggingSocketPlug{math/luamml/mtable/tag/set}{thesis}
        \NewTaggingSocketPlug{math/luamml/mtable/innertable/finalize}{thesis}
          {
            \__luamml_amsmath_finalize_inner_table:
            \bool_if:NT \l__math_gathersplit_bool { \use:c{thesis@luamml@wraptag} }
          }
        \AssignTaggingSocketPlug{math/luamml/mtable/innertable/finalize}{thesis}
      }
  }

%% ---------------------------------------------------------------------------------------------
%% 3. Theorems, proofs and other blocks that stand on their own (LaTeX's code of June 2026, which has
%% the socket block/endpe; the November 2026 code needs nothing here). That code ends them without the
%% usual note that what follows straight on continues the paragraph (\if@endpe), which its tagging
%% needs, but it changes the look: text straight after was indented, and a theorem, proof or list
%% straight after got \partopsep more space above it. Here the look stays as in the untagged editions,
%% while the tagging code still sees what it expects.
\socket_if_exist:nT { block/endpe }
  {
    \bool_new:N \g__thesis_block_endpe_bool   % a block on its own has just ended
    \bool_new:N \l__thesis_block_pre_bool
    \bool_new:N \l__thesis_block_space_bool
    \bool_new:N \l__thesis_block_after_bool
    \cs_new_protected:Npn \__thesis_block_after_standalone:
      {
        \bool_gset_true:N \g__thesis_block_endpe_bool
        \def \par
          {
            \use:c { @restorepar } \bool_gset_false:N \g__thesis_block_endpe_bool
            \everypar { } \par
          }
        \everypar
          {
            { \setbox 0 \lastbox }
            \everypar { } \bool_gset_false:N \g__thesis_block_endpe_bool
          }
      }
    \socket_new_plug:nnn { block/endpe } { thesis-standalone }
      { \legacy_if_gset_false:n { @endpe } \aftergroup \__thesis_block_after_standalone: }
    \AddToHook{begindocument}
      {
        \cs_new_eq:NN \__thesis_block_recipe_standalone: \__block_recipe_standalone:
        \cs_gset_protected:Npn \__block_recipe_standalone:
          {
            \__thesis_block_recipe_standalone:
            \socket_assign_plug:nn { block/endpe } { thesis-standalone }
          }
        \NewTaggingSocketPlug { block/recipe } { thesis }
          {
            \legacy_if:nTF { @endpe }
              { \bool_set_true:N \l__thesis_block_pre_bool }
              { \bool_set_false:N \l__thesis_block_pre_bool }
            \bool_set_eq:NN \l__thesis_block_space_bool \l__thesis_block_pre_bool
            \bool_if:NT \g__thesis_block_endpe_bool
              { \bool_set_true:N \l__thesis_block_space_bool }
            \bool_gset_false:N \g__thesis_block_endpe_bool
            \use:c { __block_recipe_#1: }
            \legacy_if:nTF { @endpe }
              { \bool_set_true:N \l__thesis_block_after_bool }
              { \bool_set_false:N \l__thesis_block_after_bool }
            \bool_if:NT \l__thesis_block_space_bool
              {
                % the space above as without tagging, then the state the tagging code set
                \legacy_if_gset_true:n { @endpe }
                \cs_set_eq:NN \__thesis_displayblock_begin: \__kernel_displayblock_begin:
                \cs_set_protected:Npn \__kernel_displayblock_begin:
                  {
                    \bool_if:NTF \l__thesis_block_after_bool
                      { \legacy_if_gset_true:n { @endpe } }
                      { \legacy_if_gset_false:n { @endpe } }
                    \__thesis_displayblock_begin:
                  }
              }
          }
        \AssignTaggingSocketPlug { block/recipe } { thesis }
      }
  }

%% ---------------------------------------------------------------------------------------------
%% 4. Figures side by side (a minipage each, with its own caption, in one figure): a figure may have only
%% one caption in a tagged PDF, so each minipage after the first starts a figure of its own in the
%% structure (the pages are unchanged).
\int_new:N \g__thesis_float_captions_int
\AddToHook{cmd/@floatboxreset/after}{ \int_gzero:N \g__thesis_float_captions_int }
\AddToHook{cmd/end@float/before}{ \int_gzero:N \g__thesis_float_captions_int }
\AddToHook{cmd/end@dblfloat/before}{ \int_gzero:N \g__thesis_float_captions_int }
\AddToHook{cmd/caption/before}{ \int_gincr:N \g__thesis_float_captions_int }
\AddToHook{env/minipage/before}
  {
    \bool_lazy_and:nnT
      { \cs_if_exist_p:c { @captype } }
      { \int_compare_p:nNn { \g__thesis_float_captions_int } > { 0 } }
      {
        \tagpdfsetup{float/split}
        \int_gzero:N \g__thesis_float_captions_int
      }
  }
\ExplSyntaxOff

%% ---------------------------------------------------------------------------------------------
%% 5. tagpdf stores the number of pages, marked-content items and structure elements of the last compile
%% (the label @tag@LastPage). The two totals change whenever anything earlier in the thesis gains or
%% loses a tagged item, such as the table of contents once the references are in it, and LaTeX would then
%% warn that labels may have changed and ask for another compile, although tagpdf (in LuaLaTeX) does not
%% use them for the PDF. Only a change in the number of pages counts here, as for any other label.
%% (latexmkrc ignores the same line.)
\ExplSyntaxOn
\cs_new:Npn \__thesis_abspage:w #1#2#3 \q_stop { \tl_if_eq:nnT {#1} {abspage} {#2} }
\AddToHook{enddocument/afterlastpage}
  {
    \cs_set_eq:NN \__thesis_label_testdef:nn \@kernel@new@label@record@testdef
    \cs_set_protected:Npn \@kernel@new@label@record@testdef #1#2
      {
        \str_if_eq:nnTF {#1} {@tag@LastPage}
          {
            \tl_if_exist:cTF { r@#1 }
              {
                \str_if_eq:eeF
                  { \exp_last_unbraced:Nv \__thesis_abspage:w { r@#1 } {} {} {} \q_stop }
                  { \__thesis_abspage:w #2 {} {} {} \q_stop }
                  { \@tempswatrue }
              }
              { \@tempswatrue }
          }
          { \__thesis_label_testdef:nn {#1} {#2} }
      }
  }
\ExplSyntaxOff

%% ---------------------------------------------------------------------------------------------
%% 6. Table cells: LaTeX gives each header cell its scope (column, row or both), and each cell that
%% spans columns its span, through attribute classes (/C /TH-col, /C /colspan-2, defined once in the
%% ClassMap). That is valid in PDF 1.7 and 2.0, and veraPDF reads it, but PAC does not: it reports
%% every header cell as having "no associated subcells", and every table with a \multicolumn as not
%% regular (tested on plain LaTeX documents in PDF 1.7 and 2.0). So the table attributes of a cell's
%% classes are also written on the cell itself (/A << /O /Table ... >>), next to the classes, when the
%% structure is written out at the end. Remove once PAC reads attribute classes, or once LaTeX writes
%% them on the cell itself (the package watch fingerprints \__tag_struct_write_obj:n, which this wraps).
\ExplSyntaxOn
\cs_if_exist:NT \__tag_struct_write_obj:n
  {
    \cs_new_eq:NN \__thesis_struct_write_obj:n \__tag_struct_write_obj:n
    \cs_set_protected:Npn \__tag_struct_write_obj:n #1
      {
        \prop_get:cnNT { g__tag_struct_#1_prop } { C } \l__thesis_tmpa_tl
          {
            % the classes: one name, or several in [ ]
            \str_set:Ne \l__thesis_tmpa_str { \tl_to_str:N \l__thesis_tmpa_tl }
            \str_remove_all:Nn \l__thesis_tmpa_str { [ }
            \str_remove_all:Nn \l__thesis_tmpa_str { ] }
            \seq_set_split:NnV \l__thesis_tmpa_seq { ~ } \l__thesis_tmpa_str
            % the entries of those owned by Table, without their /O /Table
            \str_clear:N \l__thesis_tmpb_str
            \seq_map_inline:Nn \l__thesis_tmpa_seq
              {
                \prop_get:NnNT \g__tag_attr_entries_prop {##1} \l__thesis_tmpb_tl
                  {
                    \str_set:Ne \l__thesis_tmpc_str { \tl_to_str:N \l__thesis_tmpb_tl }
                    \str_if_in:NnT \l__thesis_tmpc_str { /Table }
                      {
                        \seq_set_split:NnV \l__thesis_tmpb_seq { ~ } \l__thesis_tmpc_str
                        \seq_map_inline:Nn \l__thesis_tmpb_seq
                          {
                            \str_if_eq:nnF {####1} { /O }
                              {
                                \str_if_eq:nnF {####1} { /Table }
                                  { \tl_if_blank:nF {####1} { \str_put_right:Nn \l__thesis_tmpb_str { ~ ####1 } } }
                              }
                          }
                      }
                  }
              }
            \str_if_empty:NF \l__thesis_tmpb_str
              {
                \prop_get:cnNTF { g__tag_struct_#1_prop } { A } \l__thesis_tmpa_tl
                  { % (other attributes already: one array of them all, without nesting arrays)
                    \tl_set:Ne \l__thesis_tmpa_tl { \tl_trim_spaces:V \l__thesis_tmpa_tl }
                    \str_if_eq:eeT { \tl_head:N \l__thesis_tmpa_tl } { [ }
                      { \tl_set:Ne \l__thesis_tmpa_tl { \tl_range:Nnn \l__thesis_tmpa_tl { 2 } { -2 } } }
                    \prop_gput:cne { g__tag_struct_#1_prop } { A }
                      { [ \l__thesis_tmpa_tl \c_space_tl << /O /Table \l__thesis_tmpb_str \c_space_tl >> ] }
                  }
                  {
                    \prop_gput:cne { g__tag_struct_#1_prop } { A }
                      { << /O /Table \l__thesis_tmpb_str \c_space_tl >> }
                  }
              }
          }
        \__thesis_struct_write_obj:n {#1}
      }
  }
\tl_new:N \l__thesis_tmpa_tl
\tl_new:N \l__thesis_tmpb_tl
\str_new:N \l__thesis_tmpa_str
\str_new:N \l__thesis_tmpb_str
\str_new:N \l__thesis_tmpc_str
\seq_new:N \l__thesis_tmpa_seq
\seq_new:N \l__thesis_tmpb_seq
\ExplSyntaxOff

\makeatother
\fi
